% ==========================================
% 第 6 章：安全守则
% ==========================================
\section{实验室安全守则}

人形机器人包含高扭矩关节以及高电压、大电流的供电系统。即使机器人暂时没有移动，通电关节也可能突然施力，因此上电、测试、搬运和故障处置都需要遵循固定流程。具体风险应结合平台参数和实验室安全规程评估。

以下流程用于降低误操作和设备失控的风险。任何一项无法确认时，都应暂停上电并联系实验室负责人。

\subsection{上电之前：六项检查}

\begin{enumerate}[leftmargin=*]
    \item \textbf{吊机固定好了吗？} 保护绳受力，留 5 到 10 厘米余量。
    \item \textbf{关节处于规定位置吗？} 按平台说明确认关节姿态，避免在扭曲或机械死点状态下上电。
    \item \textbf{周围清空了吗？} 一米内没有人，地面没有线缆和工具。
    \item \textbf{急停在手边吗？} 按得到，够得着，知道按下去会怎样。
    \item \textbf{遥控器在谁手上？} 指定一个人，拇指放在阻尼键上。
    \item \textbf{三个人都在吗？} 缺一个就不上电。
\end{enumerate}

\noindent 任一项未确认时，不得上电。

\subsection{保护吊架的使用}

保护吊架用于限制机器人失衡后的下坠，不应持续牵引机器人。吊绳过紧会改变测试时的受力条件，影响步态表现。绳长、挂点和允许的松弛量应按平台说明及实验室规程设置。

\begin{itemize}[leftmargin=*]
    \item 绳子调到机器人站直时略松，留 5 到 10 厘米下沉量。
    \item 挂点在躯干重心正上方，不要在头部或者单侧肩膀。
    \item 走路测试时，吊架操作员应跟随机器人移动，并避免主动牵引机器人；吊架只承担规程规定的保护作用。
    \item 测试期间应始终保持保护绳连接。失控跌落可能损坏关节、齿轮箱或结构件。
\end{itemize}

\subsection{三人测试分工}

\begin{table}[H]
\centering
\small
\begin{tabularx}{\textwidth}{p{2.2cm} p{1.4cm} X}
\toprule
\textbf{岗位} & \textbf{代号} & \textbf{职责与站位} \\
\midrule
推吊车 & A & 双手扶吊车，视线盯着机器人重心，跟着移动。随时准备往上提。 \\
\addlinespace
遥控手 & B & 握遥控器，拇指常在阻尼键上。只盯机器人，不聊天。 \\
\addlinespace
监护 & C & 站在急停和电源旁，管人、管线、管记录。出事先拍急停。 \\
\bottomrule
\end{tabularx}
\caption{机器人测试期间三名成员的职责分工。}
\end{table}

\noindent 测试时，吊架操作员和遥控手都可能无法立即处理电源或人员隔离，因此需要独立的监护人负责急停、周边安全和异常记录。

\subsection{阻尼与急停的使用边界}

\begin{table}[H]
\centering
\small
\begin{tabularx}{\textwidth}{p{2.6cm} X X}
\toprule
& \textbf{进阻尼} & \textbf{拍急停} \\
\midrule
做了什么 & 依平台实现切换到阻尼状态，具体电气行为以厂家说明为准 & 通过硬件回路切断动力输出；机器人仍可能因重力或惯性继续运动 \\
\addlinespace
作用 & 在驱动与通信正常时，可能减少主动输出并降低部分冲击 & 可在主控制软件失效时切断动力输出，具体效果取决于硬件设计 \\
\addlinespace
限制 & 依赖控制链路和驱动器仍能响应，不适用于所有故障 & 切断动力可能使机器人失去主动支撑；应按设备设计评估后果 \\
\addlinespace
使用条件 & 按厂家说明和已验证的本队流程决定 & 存在人员危险、烟雾或火花，或其他需要立即切断动力的情况 \\
\bottomrule
\end{tabularx}
\caption{阻尼与急停的作用不同，具体处置顺序应遵循平台说明和实验室批准的流程。}
\end{table}

阻尼模式并不适用于所有故障，也不能替代硬件急停。应根据故障类型、人员位置和设备说明执行经过验证的停机流程；若存在人员受伤风险、出现烟雾或火花，或阻尼指令无响应，应立即按实验室应急流程使用硬件急停。实际处置顺序必须由 T4 厂家手册与实验室安全负责人确认。

\subsection{三类常见危险状态}

\begin{table}[H]
\centering
\small
\begin{tabularx}{\textwidth}{p{2.6cm} p{4.6cm} X}
\toprule
\textbf{类型} & \textbf{什么样子} & \textbf{为什么危险} \\
\midrule
自激振荡 & 全身高频抖动，幅度越来越大 & 控制环路相位错了，越纠越偏，几秒内能把关节甩到限位 \\
\addlinespace
单腿支撑失控 & 一条腿突然伸直或猛收 & 重心瞬间偏移，机器人朝一侧砸，几十公斤砸下来 \\
\addlinespace
线束缠绕 & 转身时线被拉住，机器人僵住 & 电机还在出力，线被扯断，或者把机器人拽倒 \\
\bottomrule
\end{tabularx}
\caption{测试过程中需要识别的三类危险状态。}
\end{table}

\noindent 出现上述现象时，应立即按实验室规定执行停机流程。不要在机器人仍可能驱动关节时靠近检查，也不要把查看日志或关闭软件作为首要处置动作。

\subsection{机器人跌倒后的处理}

\begin{enumerate}[leftmargin=*]
    \item 确认周围人员安全后，按平台规定的应急流程停机并隔离动力源。不得在关节仍可能主动施力时扶起机器人。
    \item 拍照。记录姿态和现场，事后分析靠这个。
    \item 检查三个地方：关节有没有异响、结构件有没有裂纹、线束有没有被扯。
    \item 只有在动力源已隔离、设备负责人确认可以操作后，才检查关节是否卡滞；不要强行转动可能受损的关节。
    \item 重新上电前，应先完成结构、电气和关节检查。首次测试应在保护措施到位的条件下，从低风险动作开始。
\end{enumerate}

\subsection{开机与关机顺序}

\begin{itemize}[leftmargin=*]
    \item \textbf{开机：} 低压先通（工控机、传感器），确认自检通过，喊一声「注意上高压」，再合动力电。
    \item \textbf{关机：} 按平台规定的停机顺序让机器人进入安全状态，确认机械支撑可靠后，再断开动力电并关闭工控机。
    \item \textbf{严禁}在关节扭曲、踩到异物、或者人还扶着机器人的时候合闸。
\end{itemize}

\vspace{0.4cm}
\begin{dangerbox}[贴在工位上的那张卡]
\textbf{上电前：} 吊机就位 \(\cdot\) 关节归零 \(\cdot\) 清场 \(\cdot\) 急停在手 \(\cdot\) 遥控有人 \(\cdot\) 三人到齐\\
\textbf{出事时：} 按实验室批准的应急流程停机；存在人员危险、烟雾或火花时立即使用硬件急停；隔离动力源后再接触机器\\
\textbf{原则：} 先确保人员安全，再隔离动力源并记录设备状态。
\end{dangerbox}

\newpage